\chapter{Plan de Seguridad y Salud}
\label{chap:seguridad}

La instalación del cableado implica trabajos en altura, manejo de herramientas y proximidad a instalaciones eléctricas. Este capítulo establece las medidas básicas para proteger a los técnicos y evitar daños a la infraestructura existente del pabellón.

\section{Normas de Seguridad Eléctrica y Física}
Durante la obra se respetarán las siguientes medidas de seguridad:
\begin{itemize}[label=$\bullet$,leftmargin=1.4em]
    \item \textbf{Equipo de Protección Personal (EPP):} Todo técnico usará casco, guantes, lentes de seguridad y calzado antideslizante durante la instalación.
    \item \textbf{Trabajo sin energía cercana:} Antes de intervenir cerca de tableros o tomas eléctricas, se verificará que la zona esté desenergizada o se mantendrá la separación de 150\,mm respecto a los conductores de energía.
    \item \textbf{Herramientas en buen estado:} Las herramientas de corte y ponchado se mantendrán en condiciones adecuadas para evitar cortes o accidentes.
    \item \textbf{Orden y limpieza:} Las áreas de trabajo se mantendrán libres de obstáculos y de recortes de cable que puedan causar tropiezos.
\end{itemize}

\section{Procedimientos para Trabajos en Altura y Canalizaciones}
Gran parte del tendido se realiza en techos y partes altas de los muros (puntos de Wi-Fi y CCTV), por lo que se aplican precauciones específicas:
\begin{itemize}[label=$\bullet$,leftmargin=1.4em]
    \item Uso de escaleras o andamios estables y en buen estado, nunca improvisados.
    \item Para alturas mayores a 1.8\,m, el técnico empleará arnés de seguridad anclado a un punto fijo.
    \item Una segunda persona asistirá desde el piso para estabilizar la escalera y alcanzar materiales.
    \item Al cruzar canalizaciones por techos, se evitará pisar o apoyarse sobre estructuras frágiles (cielos rasos, ductos existentes).
\end{itemize}

\section{Plan de Prevención de Riesgos}
Para anticipar los principales riesgos de la obra se elabora una matriz simple que relaciona cada peligro con su medida preventiva, como muestra el Cuadro~\ref{tab:riesgos}.

\begin{table}[H]
\centering
\caption{Matriz de identificación y prevención de riesgos.}
\label{tab:riesgos}
\begin{tabular}{lp{8cm}}
\hline
\textbf{Riesgo Identificado} & \textbf{Medida Preventiva} \\
\hline
Caídas desde altura & Uso de arnés, escaleras estables y asistente en piso. \\
%
Contacto eléctrico & Trabajar con la zona desenergizada y mantener la separación de servicios. \\
%
Cortes con herramientas & Uso de guantes y manejo correcto de las herramientas de ponchado. \\
%
Daño a instalaciones existentes & Revisar planos previos y evitar perforaciones sobre ductos o tuberías ya instaladas. \\
%
Tropiezos y caídas a nivel & Mantener orden y limpieza, recogiendo recortes de cable. \\
\hline
\end{tabular}
\end{table}

\noindent Con estas medidas se busca que la ejecución del proyecto se realice sin accidentes y sin afectar las actividades académicas que continúan en el pabellón durante la instalación.
